Difference Between Call by Value and Call by Reference in C++

Functions in C++ receive arguments either by call by value or call by reference. These parameter passing methods determine whether a function works on a copy of the data or the original variable.

Call by Value passes a copy of the argument to the function.
Call by Reference passes a reference to the original variable, allowing direct modification.

Actual Parameters vs Formal Parameters
Before understanding these methods, it is important to know the two types of function parameters.

Actual Parameters (Arguments): Values passed during the function call.
Formal Parameters: Parameters declared in the function definition that receive the arguments

For example:




#include <iostream>
using namespace std;

void display(int x) {    // x is a formal parameter
    cout << x;
}

int main() {
    int num = 10;
    display(num);        // num is an actual parameter
    return 0;
}


Call by Value
In call by value, the function receives a copy of the original variable. Changes made inside the function affect only the copied value, leaving the original variable unchanged.

Keeps the original variable unchanged after the function call.
Suitable for small data types such as int and char.
Useful when the function only needs to read the value.

#include <iostream>
using namespace std;

// Function receives a copy
void increment(int num)
{
    num++;
    cout << "Inside function: " << num << endl;
}

int main()
{
    int number = 5;

    increment(number);

    cout << "Outside function: " << number << endl;

    return 0;
}

Explanation

number is passed by value to increment(), a copy of number is stored in num.
Incrementing num changes only the copied value and the original variable number remains unchanged.


Call by Reference
In call by reference, the function receives a reference to the original variable instead of a copy. As both parameters refer to the same memory location, changes made inside the function directly modify the original variable.

Avoids copying large objects, improving performance.
Allows the function to modify the original variable.
Commonly used for large objects or output parameters.

#include <iostream>
using namespace std;

// Function receives a reference
void increment(int& num)
{
    num++;
    cout << "Inside function: " << num << endl;
}

int main()
{
    int number = 5;

    increment(number);

    cout << "Outside function: " << number << endl;

    return 0;
}

Explanation

number is passed by reference using int&, num becomes another name for number.
Incrementing num directly modifies the original variable and the updated value is visible after the function returns.

